% 原书物理页=印刷页60--63上；第4章引言与4.1全节.
\chapter{半单 Lie 代数的表示}\label{chap:ssl-4}
我们已经见过，半单 Lie 代数 $\mathfrak g$ 的有限维表示是完全可约的. 这意味着，只需分类不可约表示，就能理解 $\mathfrak g$ 的整个有限维表示理论. 不过，为了得到这一分类，我们也将考虑某些无穷维表示.
\section{最高权模}\label{sec:ssl-4-1}
设 $\mathfrak g$ 是一个半单 Lie 代数，$\mathfrak h$ 是一个 Cartan 子代数，$\Phi\subset\mathfrak h^*$ 是根集，$\Phi^\vee=\{\alpha^\vee=H_\alpha,\ \alpha\in\Phi\}\subset\mathfrak h$ 是对偶根集.
对于 $\mathfrak g$ 的一个表示 $V$ 以及 $\lambda\in\mathfrak h^*$，照常记
\begin{align*}
V_\lambda:=\{v\in V,\ \forall H\in\mathfrak h,\ H.v=\lambda(H)v\},
\end{align*}
若 $V_\lambda\neq0$，就称 $\lambda$ 为 $V$ 的一个权. 回顾：对任意根 $\alpha\in\Phi$，都有 $\mathfrak g_\alpha.V_\lambda\subset V_{\lambda+\alpha}$. 记 $P(V)$ 为 $V$ 的权集. 如果 $V=\sum_{\lambda\in P(V)}V_\lambda$，就说 $V$ 是其权子空间之和.
\begin{exercise*}
证明：若 $V$ 是其权子空间之和，则 $V$ 的任意子表示以及任意商表示也具有这一性质.（不假设 $V$ 有限维.）
\end{exercise*}
还记 $P:=\{\lambda\in\mathfrak h^*,\ \lambda(\Phi^\vee)\subset\mathbb Z\}$. 这是一个秩为 $\dim\mathfrak h$ 的自由 Abel 群，称为 $\Phi$ 的\newterm{ssl-integral-weight-lattice}{整权格}{integral weight lattice}.
\begin{proposition}\label{prop:ssl-4-1-1}
假设 $(V,r)$ 有限维. 则 $V$ 是其权子空间之和，且 $P(V)\subset P$.
\end{proposition}
\begin{proof}
我们已经见过，$\mathfrak h$ 由 $\mathfrak g$ 的半单元素组成，故 $r(\mathfrak h)$ 由 $\mathfrak{gl}(V)$ 中的可对角化元素组成. 因为 $r(\mathfrak h)$ 是一个交换子代数，这些元素可以同时对角化，第一条断言由此得到. 为证第二条，取 $\alpha\in\Phi$ 并令 $\mathfrak{sl}_\alpha=\mathfrak g_\alpha\oplus K.H_\alpha\oplus\mathfrak g_\alpha$.\footnote{译注：原文两次写作 $\mathfrak g_\alpha$，此处照录；其中一项应为 $\mathfrak g_{-\alpha}$.}
我们知道存在同构 $\mathfrak{sl}_2\xrightarrow{\sim}\mathfrak{sl}_\alpha$，把 $H$ 映为 $H_\alpha$. 通过这个同构将 $V$ 看作 $\mathfrak{sl}_2$ 的表示时，$H$ 的权是 $\lambda(H_\alpha)=\alpha^\vee(\lambda)$，其中 $\lambda\in P(V)$. 而 $\mathfrak{sl}_2$ 的表示理论告诉我们，这些权都是整数.
\end{proof}
现在固定 $\Phi$ 的一组基 $\Delta$，并记 $\Phi_+$ 为相应的正根集. 我们有分解 $\mathfrak g=\mathfrak n_-\oplus\mathfrak h\oplus\mathfrak n_+$，其中 $\mathfrak n_\pm=\sum_{\alpha\in\Phi_\pm}\mathfrak g_\alpha$.
\begin{definition}\label{def:ssl-4-1-2}
若向量 $v\in V$ 满足 $v\in V_\lambda$ 且 $\mathfrak n_+.v=0$，则称 $v$ 为权为 $\lambda$ 的\newterm{ssl-primitive-vector}{本原向量}{primitive vector}.
\end{definition}
\begin{exercise*}
证明：这等价于要求 $\mathfrak g_\alpha.v=0$ 对一切 $\alpha\in\Delta$ 成立.
\end{exercise*}
当 $V$ 有限维时，$P(V)$ 的有限性和包含关系 $\mathfrak g_\alpha.V_\lambda\subset V_{\lambda+\alpha}$ 保证了 $V$ 中存在本原向量. 若 $V$ 还是不可约的，则可见 $V$ 由一个本原向量生成. 这就引出了``最高权''表示的概念.
\begin{definition}\label{def:ssl-4-1-3}
若一个表示 $V$（未必有限维）由一个本原向量 $v$ 生成，则称它为\newterm{ssl-highest-weight-module}{最高权模}{highest weight module}. 此时称 $v$ 为一个``最高权向量''，并称 $v$ 的权 $\lambda$ 为``$V$ 的一个最高权''.
\end{definition}
这一术语可由下面结果的第 \ref{ssl-4-1-4-i} 项得到解释.
\begin{proposition}\label{prop:ssl-4-1-4}
设 $V$ 是一个最高权为 $\lambda$ 的表示.
\begin{enumerate}[label=\roman*)]
\item\label{ssl-4-1-4-i} $P(V)\subset\lambda-\sum_{\alpha\in\Delta}\mathbb N\alpha$. 换言之，对于 $\mathfrak h^*$ 上由 $\lambda'\preccurlyeq\lambda\Longleftrightarrow\lambda-\lambda'\in\sum_{\alpha\in\Delta}\mathbb N\alpha$ 定义的偏序，$\lambda$ 是 $P(V)$ 的最大元.
\item $V=\sum_{\mu\in P(V)}V_\mu$，即 $V$ 是其权子空间之和.
\item 每个 $V_\mu$ 都有限维，且 $\dim V_\lambda=1$.
\end{enumerate}
\end{proposition}
\begin{proof}
令 $\mathfrak b:=\mathfrak h\oplus\mathfrak n_+$. 把表示 $V$ 延拓到泛包络代数 $\mathcal U\mathfrak g$ 后，直线 $Kv$ 在 $\mathcal U\mathfrak b$ 下稳定. 另一方面，若对每个 $\alpha\in\Phi_+$ 选一个 $\mathfrak g_{-\alpha}$ 中的非零元素 $F_\alpha$，则 Poincaré--Birkhoff--Witt 定理告诉我们，$\mathcal U\mathfrak g$ 作为右 $\mathcal U\mathfrak b$-模，由形如 $F_{\alpha_s}F_{\alpha_{s-1}}\cdots F_{\alpha_1}$ 的元素生成，其中 $s\in\mathbb N$、$(\alpha_i)_i\in(\Phi_+)^s$（若 $s=0$，约定这个乘积为 $1$）. 换言之，
\begin{align*}
\mathcal U\mathfrak g=\sum_{s\in\mathbb N}\sum_{(\alpha_i)_i\in(\Phi_+)^s}F_{\alpha_s}F_{\alpha_{s-1}}\cdots F_{\alpha_1}\mathcal U\mathfrak b.
\end{align*}
由此得到，$V$ 作为 $K$-向量空间，由形如 $F_{\alpha_s}F_{\alpha_2}\cdots F_{\alpha_1}.v$ 的向量生成.\footnote{译注：原文此处乘积的第二个因子写作 $F_{\alpha_2}$，与上一行的 $F_{\alpha_{s-1}}$ 不同，此处照录.}
这样的向量的权为 $\lambda-\sum_{i=1}^s\alpha_i$. 于是得到前两项以及 $\dim(V_\lambda)=1$. 为了说明 $\dim V_\mu$ 有限，只需注意，满足 $\mu=\lambda-\sum_{i=1}^s\alpha_i$ 的序列 $\alpha_1,\alpha_2,\ldots,\alpha_s$ 只有有限多个.
\end{proof}
\begin{remark*}
最高权表示 $V$ 是不可分解的. 事实上，若 $V=V_1\oplus V_2$ 是一个在 $\mathfrak g$ 下稳定的分解，则 $V_\lambda=V_{1,\lambda}\oplus V_{2,\lambda}$. 但因为 $\dim(V_\lambda)=1$，存在 $i=1$ 或 $2$ 使得 $V_\lambda=V_{i,\lambda}$. 而 $V$ 由 $V_\lambda$ 生成，故 $V\subset V_i$.
\end{remark*}
\begin{numpara}[Verma 模]\label{np:ssl-4-1-5}
上面的证明启发我们定义一个``泛''最高权表示. 令
\begin{align*}
M(\lambda):=\mathcal U\mathfrak g\otimes_{\mathcal U\mathfrak b}K_\lambda.
\end{align*}
这是由线性形式 $\mathfrak b\twoheadrightarrow\mathfrak h\xrightarrow{\lambda}K$ 给出的一维 $\mathcal U\mathfrak b$-模所``诱导''的 $\mathcal U\mathfrak g$-模. 由张量积的伴随性质，对任意表示 $V$ 都有
\begin{align*}
\operatorname{Hom}_{\mathfrak g}(M(\lambda),V)\simeq\operatorname{Hom}_{\mathfrak b}(K_\lambda,V)=\{v\in V_\lambda,\ \mathfrak n_+v=0\}.
\end{align*}
特别地可见，每个最高权为 $\lambda$ 的表示都是 $M(\lambda)$ 的商. 后者称为``权为 $\lambda$ 的\newterm{ssl-verma-module}{Verma 模}{Verma module}''.
\end{numpara}
\begin{theorem}\label{thm:ssl-4-1-6}
（这里的表示可以是无穷维的.）
\begin{enumerate}[label=\roman*)]
\item\label{ssl-4-1-6-i} 对任意 $\lambda\in\mathfrak h^*$，存在最高权为 $\lambda$ 的不可约表示 $V(\lambda)$，且在同构意义下唯一.
\item\label{ssl-4-1-6-ii} 若 $V(\lambda)\simeq V(\lambda')$，则 $\lambda=\lambda'$.
\end{enumerate}
\end{theorem}
\begin{proof}
先证第二项. 设 $V(\lambda)\simeq V(\lambda')=:V$. 则 $P(V)\subset\lambda-\sum_{\alpha\in\Delta}\mathbb N\alpha$，故 $\lambda'\preccurlyeq\lambda$. 但也有 $P(V)\subset\lambda'-\sum_{\alpha\in\Delta}\mathbb N\alpha$，故 $\lambda\preccurlyeq\lambda'$. 于是 $\lambda=\lambda'$.
再证第一项. 设 $V$ 是一个最高权为 $\lambda$ 的表示，例如 $V=M(\lambda)$. 设 $W_1,W_2$ 是 $V$ 的两个\newterm{ssl-proper-submodule}{真子模}{proper submodule}. 由上述命题，$V$ 是其权子空间之和，故由上面的一道习题，每个 $W_i$ 也是其权子空间之和. 因为 $W_i$ 是真子模，所以 $W_i\subset\sum_{\mu\prec\lambda}V_\mu$，从而和 $W=W_1+W_2$ 也包含在 $\sum_{\mu\prec\lambda}V_\mu$ 中，因此也是 $V$ 的一个真 $\mathfrak g$-子模. 所以可以考虑 $V$ 的最大真子模 $U$. 明确地说，$U=\sum_{W\subsetneq V}W$，并且 $U\subsetneq V$. 于是商 $\overline V:=V/U$ 非零，最高权为 $\lambda$（以 $V$ 的一个本原向量的像为本原向量），而且不可约，因为若 $\overline W\subset\overline V$ 非零，则 $\overline W$ 在 $V$ 中的原像 $W$ 严格包含 $U$，故等于 $V$. 此外，这一构造表明，$\overline V$ 是 $V$ 唯一的\newterm{ssl-irreducible-quotient}{不可约商}{irreducible quotient}（因为投影 $V\twoheadrightarrow\overline V$ 通过任意非零满态射 $V\twoheadrightarrow V'$ 分解）.
剩下要说明，最高权为 $\lambda$ 的所有单模都同构. 但正如定理前所见，这样的模都是 $M(\lambda)$ 的商，故同构于上面构造的 $M(\lambda)$ 的唯一不可约商 $V(\lambda)$.
\end{proof}
